\section{Discussion}
\label{sec:discussion}

This section interprets the results rather than restating them.  The discussion covers what the partition means, which sector of the model provides the branch, how the
mechanism relates to published formulations, where each constant comes from and
what the model does not resolve.

\subsection{The observable: a modal partition of the scattered power}
\label{sec:disc_meaning}

The central quantity of this work is $R_B$: the share of the incident
power that ends up in one particular outgoing mode, measured at a plane ten
thicknesses from the edge.  It is a partition of the scattered field and not a reflectance.  At the operating point the backward mode takes
\SI{1.93}{\percent} while the remaining propagating outgoing modes take
\SI{98.05}{\percent}; at the maximum the backward mode takes
\SI{94.29}{\percent} while the rest take \SI{5.76}{\percent}.  At the operating
point the total outgoing propagating power is $0.99980$ of the incident power
and at the maximum $1.00050$, so the plate is an almost perfect scatterer of the
guided mode in both cases and the question of interest is not how much is reflected but into which branch it goes.

The distinction has a practical consequence.  A measurement of the total
reflectance of an edge would see essentially unity everywhere on the branch.  A
measurement that resolves the modal content (by a wavenumber filter, or by a
time-of-flight separation of the branches) would see the partition move from
almost nothing in the backward channel to more than half of the incident power.
The negative reflection discussed here is an observable that requires modal
resolution.

The word ``backward'' is used in an operational sense: a mode is backward when
the $x$ component of its thickness-integrated Poynting flux is negative, and
\eqref{eq:classification} confirms that the same mode is outgoing at the edge.
The dispersion branch has a region where $v_g$ and $v_E$ are negative
(\cref{fig:fig01}b), which is necessary for a backward mode to exist, and that
region is a property of the infinite plate.  What creates the conversion is the
edge, and what controls its strength is where the operating point lies on the branch.

\subsection{The mechanism}
\label{sec:disc_mechanism}

The sweep along the branch (\cref{fig:fig04}) places the maximum of the
conversion at $k_xd=0.52719$, where the group velocity vanishes
($+0.01\,\si{\metre\per\second}$ against $-187$ and $+409\,\si{\metre\per\second}$
at the ends of the solved window) and where the backward wavenumber approaches
the incident one to within \SI{0.4}{\percent}.  This is a slow-mode effect and
it has a simple reading: at a zero-group-velocity fold the incident and backward
wavenumbers are nearly equal, so the two modes are nearly parallel and the edge
condition can be met with only a small admixture of the other outgoing modes;
the projection onto the backward mode then takes almost all of the partition.
Panel (c) of the same figure makes this explicit: all coupling levels collapse
onto one curve in $k_{Bx}/k_{ix}$, with the conversion maximal as the ratio
approaches unity.  The same mechanism explains why the conversion falls steeply
on both flanks: on the fast side the wavenumbers separate, and on the far side of
the fold the branch has changed character.  A grid of $\Delta k_xd=0.05$ would miss this maximum, because the entire top of the response is $10^{-3}$ wide; the two rasters of \cref{fig:fig04} locate it.

The reading has a second half.  At the fold the incident
mode itself is nearly cut off: its group velocity vanishes, so the flux it
carries per unit normalization falls by a factor $131$ with respect to the
operating point.  The partition is maximal at the very point where the transported power is minimal.  A statement that ninety-four percent of the energy is turned around
would be true as a ratio and false as an energy budget, and only the first is claimed.

The reflection is present before the coupling is switched on.  The
controls of \cref{sec:results_mechanism} establish this directly: with $H=0$ the
same branch converts \SI{94.29}{\percent} of the incident power at the
maximum and \SI{1.93}{\percent} at the operating point, and at the same
frequency the classical plate of the control row has no backward member.  The strong negative reflection is thus a property of the backward guided branch and of the edge,
not a resonance of the added coupling; the gradient length places the fold of
that branch at the reported frequency.

However, the coupling $H$ acts on the loss, and through the loss on the partition.  It
shifts the dispersion by at most $1.1\times10^{-5}$ in relative terms
(inset of \cref{fig:fig01}a), it shifts the loss by up to a factor of sixty-two
(\cref{fig:fig06}c,d), and the partition follows with a weight that depends on
where the branch is sampled: \SI{-1.21}{\percent} at the operating point,
\SI{-10.17}{\percent} at $k_xd=0.55$ and \SI{-36.1}{\percent} at the maximum.  The ordering (dispersion unaffected,
partition affected only where the branch is slow, loss strongly affected
everywhere) is what a carrier-level coupling acting through conduction loss
should produce.  The third of the partition lost at the maximum is mostly
loss acting on the backward mode between the edge and the reference plane:
referred to the edge plane, where both effects act at the same place, the same
comparison gives \SI{13.8}{\percent}, against \SI{36.1}{\percent} at the reference plane, so
22.3 percentage points are propagation loss and the rest is an actual reduction of the conversion.  Re-searching the extremum at each coupling level separates the
two effects further: the location moves from $0.5272$ to $0.5222$ and $0.5240$
at the two extremes of $H$, and the largest attainable share falls from $0.9429$
to $0.6034$ and $0.6832$, so the coupling both lowers the attainable maximum and displaces it, asymmetrically in the sign of $H$.

The transverse wavenumber is the organizing parameter at the operating point,
not the coupling.  \Cref{fig:fig07} shows that $k_y$ changes the backward
channel by a factor of $1.41$ across the solved box while $H$ moves it by
\SI{1.21}{\percent} branch-consistently (\SI{2.5}{\percent} at fixed frequency).
At the maximum, however, the coupling moves the partition by \SI{36.1}{\percent}, more than the
$k_y$ factor found at the operating point; the $k_y$ sweep was not repeated there.  The ordering is not a property of the coupling alone; it is a
property of the coupling and the operating point together.
Any statement about the added coupling must be made at fixed $k_y$,
and any experimental attempt to see the coupling must control the transverse
wavenumber much better than it controls the doping or the mobility.

\subsection{Relation to the literature}
\label{sec:disc_literature}

The formulation belongs to the flexoelectric-semiconductor continuum developed
by Qu and co-workers \cite{qu2020mobile,qu2022thermal}, in which the strain
gradient polarizes the medium, the polarization is screened by the mobile
carriers, and the redistributed carriers act back on the field.  The wave theory used here is that of Yang et al.\ \cite{yang2018lamb,yang2020rayleigh}
for flexoelectric plates and half-spaces with strain-gradient elasticity, and Zhang et al.\ \cite{zhang2025pscplate} find, for a simply supported plate under static load and in free vibration, that the semiconductor coupling is active only within a window of initial electron concentration, and that this window moves to larger concentrations as the plate becomes thinner.  The concentration at which their damping is largest is $10^{18.6}$, $10^{19.6}$ and $10^{20.6}\,\si{\per\cubic\metre}$ for $h=30$, $3$ and $0.3\,\si{\micro\metre}$.  The present plate ($d=\SI{200}{\nano\metre}$, $n_0\le3\times10^{18}\,\si{\per\cubic\metre}$) lies below these values, so the carrier results here describe the weakly conducting side of that window.  Qu et al.\ \cite{qu2023drag} derive plate equations for flexoelectric semiconductors by the Mindlin method with the same simplified gradient modulus $g=\ell^2c$, and check them against three-dimensional dispersion relations for a silicon plate of thickness \SI{0.2}{\micro\metre}.  For piezoelectric-semiconductor ZnO plates, Li et al.\ \cite{li2023sh} report for shear-horizontal waves that an initial density below $2\times10^{18}\,\si{\per\cubic\metre}$ barely changes the cut-off frequencies and that above $10^{23}\,\si{\per\cubic\metre}$ the plate responds as an elastic one; the density sweep of \cref{sec:results_sens} lies within this range.  Dispersion and damping of Rayleigh-type waves in a piezo-thermo-electric semiconductor with fractional and nonlocal effects have also been computed \cite{aldandani2026kg}.

What is different here is the object of study.  The references above compute
dispersion relations and attenuation of modes in an unbounded or layered
medium, and the scattering calculations available for flexoelectric media are
plane-wave reflection coefficients at a free surface of a half-space
\cite{mondal2022scattering}.  This paper computes the \emph{scattering} of a
guided mode at a free edge of a plate, resolves the outgoing spectrum into its
branches, and reports the partition.  That is a different calculation with a different output: a
dispersion curve is an eigenvalue problem, whereas an edge partition requires
the matching of a continuum of outgoing modes and a classification of each of
them.

The closest primary work is the thermodynamically consistent theory of Qu, Pan,
Zhu and Deng \cite{qu2025thermodynamic}, who derive the constitutive equations
for a flexoelectric semiconductor from the free-energy imbalance with the
Coleman--Noll procedure, and who obtain the drift--diffusion current, the
flexoelectric and piezoelectric couplings and the higher-order stress in one
framework.  That article builds a continuum theory for flexoelectric semiconductors from
the balances of mass, charge, momentum and angular momentum together with the
thermodynamic laws.  Its constitutive equations follow from the dissipation
inequality and include piezoelectric, flexoelectric, thermoelectric and
drift--diffusion couplings, and the theory is completed by a principle of virtual
work and a linearization for small carrier perturbations.  The interaction of
strain gradient and electric current is therefore the subject of that framework,
and a carrier--strain-gradient coupling of the class used here is not claimed as
new to it.  The explicit form of its free energy is reproduced, with attribution,
in the open-access computational paper of Hu et al.\ \cite{hu2025computational}
(Eqs.~(4a)--(4e) there), and that form is used for the comparison.  In
that formulation the electric Gibbs free energy is additive,
\begin{equation}
\Psi_{\text{Gib}}=\Psi_{\text{Mech}}(\varepsilon,g)
+\Psi_{\text{Elec}}(e)+\Psi_{\text{Diele}}(\varepsilon,g,e)
+\Psi_{\text{Carrier}}(n,p),
\label{eq:qu_decomposition}
\end{equation}
with
$\Psi_{\text{Mech}}=\tfrac12C_{ijkl}\varepsilon_{ij}\varepsilon_{kl}
+\tfrac12A_{ijklmn}g_{ijk}g_{lmn}$,
$\Psi_{\text{Diele}}=-\eta_{ijk}e_i\varepsilon_{jk}
-\mu_{lijk}e_l g_{ijk}$ and
$\Psi_{\text{Carrier}}(n,p)$ a function of the two carrier densities alone.
Three consequences follow for the present paper.

First, the carrier energy in \eqref{eq:qu_decomposition}, which is the form of the main text of \cite{qu2025thermodynamic} (Eqs.~(3.16)--(3.17) there), contains no strain
gradient.  Since the thermodynamic restriction fixes
$\partial\Psi_{\text{Gib}}/\partial n=\mu_n+q\psi$, the electrochemical potential
of that decomposition contains no strain-gradient contribution, and the strain
gradient reaches the carriers only \emph{indirectly}, through the flexoelectric
polarization that enters Gauss's law.  The direct bilinear term
$\Delta n\,H_{jkl}\varepsilon_{jk,l}$ of \eqref{eq:added_coupling} is absent from
that postulate, but not from the article: its Appendix~C.1 adds a chemo-mechanical
energy that is linear in the carrier concentration and in the strain and the strain gradient,
$\tilde n\,M^{n}_{IJ}G_{IJ}+\tilde n\,M^{n}_{IJK}\eta_{IJK}$ (and the same for holes), so that the
electrochemical potential acquires $M^{n}_{IJ}G_{IJ}+M^{n}_{IJK}\eta_{IJK}$ (their Eqs.~(C.1)--(C.2)).  The term used here
corresponds to $M^{n}_{IJK}$, with $H_{jkl}$ linearized about $n_0$ and written in the scaled units; the article neglects it in its Sec.~3.4 and does not pose an edge-scattering problem.
Relative to the main-text postulate it is an \emph{independent} addition, neither a
change of variables nor a special case of a more general tensor.

Second, the symbol $H$ is overloaded in the literature and we flag the collision
explicitly to avoid a misreading: in \cite{hu2025computational} the sixth-order strain-gradient elasticity tensor is written $H_{ijklmn}$ and in
\cite{qu2025thermodynamic} it is $G_{IJKLMN}$, while there $H_{IJKLM}$ is a fifth-order strain--strain-gradient
coupling; in this paper the sixth-order tensor is written $A_{ijklmn}$ (\cref{eq:enthalpy}, \cref{eq:tau});
the coefficient $H_{jkl}$ used here is a third-order carrier--strain-gradient
coefficient with different units (J\,m).  The two are unrelated, and the
positivity bounds derived in the formulation apply only to the latter.

Third, the comparison locates the contribution.  Were the coupling the main result, the contribution would be a constitutive hypothesis of the kind that \cite{qu2025thermodynamic} already systematizes; it is not the main result, because \cref{sec:results_mechanism} shows that the negative reflection exists with $H=0$.  What this paper contributes is the \emph{scattering} problem: the
partition of a guided mode into the outgoing branches of a free edge, with a
backward branch on which energy and phase travel in opposite directions, together with the demonstration that a carrier--strain-gradient enthalpy can be
added to a strain-gradient piezoelectric semiconductor without violating the
positivity bound, and a quantified account of its incremental effect.

For this reason the added term is not put forward as new.  The formulations
compared above already contain a flexoelectric coupling of strain gradients to the
electric field, and the carrier density enters them through the free energy of
the electron gas; the term $\Delta n\,H_{jkl}\varepsilon_{jk,l}$ with its
reciprocal partners in $\sigma_{ij}$ and $\widetilde\mu$ is written here as the
minimal extension that closes the enthalpy, and the paper establishes as much
about what it does not do as about what it does: it moves the partition by
\SI{1.21}{\percent} at the operating point, \SI{10.17}{\percent} at
$k_xd=0.55$ and \SI{36.1}{\percent} at the maximum, all at $|H|=0.05H^{*}$, and
it is not the origin of the backward branch.  No priority is claimed for it, as the term has precedent in the cited thermodynamic framework.  What the
present numbers do establish is that the edge-scattering problem is well posed,
that its energy budget closes to better than $10^{-11}$, and that the coupling is a perturbation of a branch that the classical plate already has, near $k_xd=0.82$, while at the operating frequency and in-plane wavenumber considered here the gradient-free plate has no accessible outgoing backward channel.

\subsection{Material provenance}
\label{sec:disc_material}

The provenance of the parameters is mixed and is stated explicitly in
\cref{tab:provenance}.  The classical constants are literature values; the
transport parameters are representative values for n-type ZnO at room
temperature linked by the Einstein relation; the strain-gradient length, the
flexoelectric coefficient and the coupling scale are \emph{assumed}.  The flexoelectric coefficient lies in the nC/m range reported for ZnO \cite{hu2025computational,zhao2021effects}.  For $H_{jkl}$ no published value is available, and the numbers used here
($|H|\le\SI{2.0e-23}{\joule\metre}$) are fixed by the positivity bound
\eqref{eq:schur}, which leaves the $(\varepsilon,\eta,\Delta n)$ block of the enthalpy positive definite with a margin of a
factor of forty-five in $H$; no specimen provides them.

\begin{table}[t]
\centering
\caption{Provenance of each parameter.}
\label{tab:provenance}
\small
\begin{tabular}{@{}p{0.2\textwidth}p{0.42\textwidth}p{0.3\textwidth}@{}}
\toprule
Parameter & Value and provenance & Source \\
\midrule
$c_{ij}$, $e_{ij}$, $\kappa_{ij}$, $\rho$ &
$c_{11}=210$, $c_{33}=221$, $c_{44}=43\,\si{\giga\pascal}$;
$e_{33}=1.32$, $e_{31}=-0.57$, $e_{15}=-0.48\,\si{\coulomb\per\square\metre}$;
$\kappa_{33}=8.85\times10^{-11}\,\si{\farad\per\metre}$; literature values &
Table~1 of \cite{zhang2025pscplate} \\
$\mu_e$, $D$ & $0.0175\,\si{\square\metre\per\volt\per\second}$,
$D=4.55\times10^{-4}\,\si{\square\metre\per\second}$, equal to $\mu_e k_BT/q$ at
$T=\SI{301.7}{\kelvin}$ (\SI{0.6}{\percent} above the value at
\SI{300}{\kelvin}); representative, within the $120$--$\SI{205}{\square\centi\metre\per\volt\per\second}$ compiled for nominally undoped bulk and film ZnO in Table~VII of \cite{ozgur2005zno}; \cite{zhang2025pscplate} use $100\,\si{\square\centi\metre\per\volt\per\second}$ and $2.6\,\si{\square\centi\metre\per\second}$; \cite{li2023sh} use $1\,\si{\square\metre\per\volt\per\second}$ and $0.026\,\si{\square\metre\per\second}$ &
input value \\
$n_0$ & $1\text{--}3\times10^{18}\,\si{\per\cubic\metre}$; modeling scenario, below the $8\times10^{13}$--$5\times10^{17}\,\si{\per\cubic\centi\metre}$ ($8\times10^{19}$--$5\times10^{23}\,\si{\per\cubic\metre}$) of Table~VII of \cite{ozgur2005zno} &
parameter sweep of \cref{fig:fig02} \\
$d$ & \SI{200}{\nano\metre} reference, 120--\SI{320}{\nano\metre} sweep;
modeling scenario & geometry sweep of \cref{fig:fig03} \\
$\ell$ & \SI{50}{\nano\metre}; assumed; the same $g=\ell^2c$ form is used in \cite{qu2023drag} with \SI{20}{\nano\metre} for silicon & this work \\
$f_{3zz}$ & \SI{1}{\nano\coulomb\per\metre}; assumed, in the nC/m range reported for ZnO; contributes
\SI{-0.01}{\percent} to $R_B$ & this work \\
$H_{jkl}$ & $H_{333}$ only, $|H|\le2.0\times10^{-23}\,\si{\joule\metre}
=0.022\,H_{\text{bound}}$; assumed, within the positivity bound & this work \\
\bottomrule
\end{tabular}
\end{table}

\subsection{Limitations}
\label{sec:disc_limitations}

Five limitations bound the conclusions, and each is stated where it enters the
model.

\begin{enumerate}
\item \emph{Geometry.}  The formulation is two-and-a-half-dimensional: a plate
infinite in $y$, semi-infinite in $x$, with three displacement components on the
$(x,z)$ cross-section.  It has one edge and not two, and it does not resolve the
three-dimensional corner where the edge meets the free surfaces.  A finite
plate, a second edge, and the interference between them are outside the model.
\item \emph{The corner.}  The corner line force \eqref{eq:corner_force} is a
natural datum that the polynomial basis cannot annihilate pointwise, so it is
measured instead of assumed, on the coherent sum of the incident mode and every
outgoing mode.  Relative to the incident mode's own corner scale it is
$8.4\times10^{-2}$ at the operating point, $3.3\times10^{-3}$ at the corrected
maximum and $3.8\times10^{-2}$ at $k_xd=0.55$, all at the lowest degree tested
(10), and it falls to $4.5\times10^{-2}$, $1.7\times10^{-3}$ and
$2.0\times10^{-2}$ at degree 28.  Over that range it follows one algebraic law in
each configuration, $3.5\times10^{-1}n^{-0.62}$,
$1.4\times10^{-2}n^{-0.63}$ and $1.6\times10^{-1}n^{-0.63}$: the same exponent to
$0.014$ in three independent configurations, which is the signature of an unresolved corner singularity of the continuum solution and not of a quadrature or coding artifact; it is unchanged to six digits when the quadrature rule is varied from 96 to 192 nodes.  So two things are established and one is not.  Established: the line force decays towards zero
algebraically and is bounded by the numbers above over the full tested range;
and it does not decide the result, because the reported partition at the maximum moves
by $3.9\times10^{-6}$ between degrees 10 and 28, while the line force halves over the same range and the energy identity closes to better than $10^{-11}$.  Not
established: that the force vanishes, or that the exponent is $0.62$ as opposed
to a nearby value; a proof would need the local corner eigenvalue problem of the
second-gradient operator with the flexoelectric and carrier couplings.  The closure claimed here is thus a \emph{bounded} one, not a full one, and the
scattering result does not rest on it.  The same holds for the traction itself: the weak
form of the edge condition closes to machine precision, but the pointwise traction is not identically zero.  At the
operating point its normalized norm away from the corners is $4.3\times10^{-2}$ at degree 10 and
$3.6\times10^{-2}$ at degree 28, while its maximum over the edge, which sits at the corners,
grows from $0.77$ to $2.4$ over the same range.  The free edge is therefore satisfied in the weak sense,
and a direct two-dimensional $x$--$z$ finite-element solution of the same variational problem, with the edge cells refined down to $0.00125\,d$, reproduces the modal backward share at the operating point to within $0.5$~\% (\cref{tab:fe2d}), with no systematic drift under refinement.  That check is made at the operating point only.  The same measurements settle what the
earlier band-energy argument could not: the strain-gradient energy in the outer
band $|z|>0.4\,d/2$ is $0.93$--$1.00$~\% of the total at the operating point,
$1.22$--$1.30$~\% at the maximum and $1.18$--$1.26$~\% at $k_xd=0.55$, and its
share decreases with degree, but the band average of the mixed hyperstress
exceeds the pointwise corner value by factors of five to ten, so no band average
is used here to argue the corner datum away.
\item \emph{The carrier model and the electrical condition.}  The only loss in the base model
is carrier conduction and diffusion; \cref{sec:results_sens} shows that an added
elastic loss tangent of $10^{-5}$ already lowers the maximum conversion from $0.943$ to $0.600$.  The carrier
physics is linearized about a uniform doping and zero bias, so the nonlinear
drift--diffusion regime, in which the perturbation is not small compared with
$n_0$, is outside the model.  The open-circuit condition is imposed weakly on the full potential space; it
is checked against the independent $C^{1}$ discretization, but it belongs to the same model and therefore serves as a consistency check of the discretization and not as a check of the physics.
\item \emph{The parameters.}  The classical constants are those of a bulk
crystal and the transport parameters are representative values for n-type ZnO at
room temperature; the strain-gradient length, the flexoelectric coefficient and
the coupling scale are assumed, and no published value of $H_{jkl}$ is available.
The reported percentages should be read as the response of this model
and not as predictions for a particular film, and the flexoelectric sector, which
is present but subdominant, is not a calibrated
mechanism.
\item \emph{The constitutive attribution.}  The relation of the added coupling
to the published flexoelectric--semiconductor formulations is discussed in \cref{sec:disc_literature}, and no priority is claimed for it.  What the
model does establish is that a bilinear carrier--strain-gradient enthalpy can be
added inside the positivity bound, that its reciprocal partners follow from the
same potential, and that its incremental effect on the partition is measurable,
of one sign, and small compared with the effect of the branch itself.
\end{enumerate}
